% SCRUM-117. Canonical sources: deployment.md and delivery dossier, audit 13/11/2026.
\section{Chiến lược và kết quả kiểm thử}
\label{sec:a2-implementation-tests}
\subsection{Mục tiêu, mức kiểm thử và môi trường}
Chiến lược kiểm thử tập trung vào rủi ro tài chính và toàn vẹn: tính giá,
ownership, idempotency, stock/coupon concurrency, payment finalization và
callback trùng/trễ. Unit test kiểm hàm/policy có cô lập; integration test kiểm
tương tác service/API và transaction; E2E HTTP kiểm boundary request/response.
E2E HTTP/JSDOM không tự thay browser end-to-end menu--confirmation hoặc giao
dịch thật với Sandbox.

API có scripts test:unit, test:integration, test:e2e tại apps/api; web có
test:unit/test tại apps/web. CI Node24 chạy npm ci, lint, unit và build, API
generate Prisma trước test; Docker dùng Node26. PostgreSQL opt-in phải dùng
DB A2 cô lập theo hướng dẫn test, không production/A1. Seed có thể tái tạo fixture nhưng
không chứng minh fresh-database setup đã được kiểm thử.

Test inventory gồm auth/products/orders, inventory PostgreSQL, coupon,
payment/idempotency/finalization concurrency, IPN--QueryDR và reconciliation;
web có checkout/payment/history/confirmation UI suites. CI web test:unit
không chứa payment-confirmation-flow.test.ts dù script test có file đó.
Vì vậy tên Unit tests hoặc CI xanh không có nghĩa mọi test trong repo đều chạy.

\subsection{Kết quả kiểm thử kế thừa theo source}
Nguồn SCRUM94 trong archive, được PR143 đưa vào tag A1-final, ghi các lượt
chạy06--07/11/2026 (Asia/Bangkok). Runtime/test source là
\path{fb293e2ac35c0fe29fbe3573bf5edf81b7a586a0}, không phải SHA tag1951dc1.
Môi trường host Windows/Node22.18/npm10.9.3, Docker Desktop4.97/Engine29.9.2/
Compose5.6.1; test container Node26.13 và PostgreSQL18.

API unit300 PASS; API E2E227 PASS; selected integration147 passed/10 skipped;
concurrency83 PASS là lượt riêng, không cộng trùng suite. Frontend174/174
PASS trong Docker DB-guard rerun; lượt host trước đó có172 passed/2 cancelled
và exit1; kết quả rerun được ghi riêng với điều kiện thực thi của nó.

Kiểm coupon/migrate tại \texttt{eef344c24fedd7480d06c085da3fff298def90d4}
khác lượt runtime nói trên. Compose verifier post-sync tại checkout
\texttt{304a7c6}, code/dev merge \texttt{dba36a5c31ea07e88b2201e3c8ffd6b345720afa},
chạy isolated test profile với VNPay config synthetic, không gọi gateway.
Default Compose từng bị chặn thiếu VNPAY\_TMN\_CODE; verifier PASS sau đó
không chứng minh default/Sandbox đã hết lỗi. Các số liệu này là evidence
kế thừa A1/Phase2, không phải kết quả chạy lại trên phiên bản mới.

PR145 báo manual regression catalog9 sản phẩm, Auth7/7 và Order/VNPay Sandbox/
History đạt trong phạm vi PR; QueryDR Sandbox E2E vẫn deferred, payment-status
route còn gap và chưa public production tại mốc A1. ADR về IPN Sandbox có
phạm vi/ngày riêng, không thay kiểm thử QueryDR lost-IPN.
Nguồn đầy đủ: \path{docs/a1-archive/verification/SCRUM94_Verification_Evidence.md}
và PR143/145 của repository A1 \cite{a1repo}.

\subsection{Kết quả MVP Sprint1 và giao diện lịch sử}
Report SCRUM-112 kiểm source
\path{34cbbc088117d7a4908d07a66cfbd03b0c905dea}, report vòng1
\path{ffbaeeb3afece5d491e25e494b099496e8234ccd}.
Report ghi API unit325, Products HTTP E2E11 và web166, tổng502/502 PASS;
Prisma validate/build PASS. API lint1 warning, web23 warnings.
Môi trường browser local là localhost:3000/API127.0.0.1:3001,
viewport320x640,375x812,768x1024; report không ghi thời điểm chạy và OS riêng.
Các ảnh ở hình\ref{fig:historical-ui} thuộc report này.

Báo cáo kiểm thử ghi nhận sai khác rounding dù các bài kiểm thử đã chạy đạt. Kết quả giao diện dùng Next.js16.3.8, khác bản16.4.6 trong source hiện hành; vì vậy số liệu được giữ theo môi trường gốc.

\subsection{Tổng hợp kiểm thử tích hợp}
\begin{longtable}{>{\raggedright\arraybackslash}p{0.24\linewidth}p{0.66\linewidth}}
\caption{Kết quả kiểm thử theo phiên bản và môi trường}\label{tab:test-results}\\
\toprule Nguồn và môi trường & Kết quả và phạm vi \\ \midrule\endfirsthead
\toprule Nguồn và môi trường & Kết quả và phạm vi (tiếp) \\ \midrule\endhead
PR25, source \texttt{2216b7a}, báo cáo ngày13/11/2026 & Bộ focused tests 38/38; Vitest 479 đạt, 6 bài PostgreSQL opt-in bỏ qua; TypeScript và build đạt. Nguồn chưa ghi thời điểm và môi trường riêng của mỗi lượt; các bài bỏ qua không được tính đạt. \\
PR26, backend Render source \texttt{2216b7a}, frontend Vercel & Báo cáo giao dịch Sandbox cho thấy Order PAID, Payment SUCCESS và loyalty +10 qua Order History. Kết quả giới hạn ở kịch bản được báo cáo; không suy thành kiểm thử độc lập mọi nhánh IPN hoặc QueryDR. \\
CI ứng dụng, source \texttt{0e43bcb}, 13/11/2026 & API và Web hoàn tất install, lint, unit test và build trên Node24. Pipeline không chạy browser E2E, migration production hoặc build báo cáo LaTeX. \\
HTTP public, 14/11/2026, 10:31 UTC & Frontend, proxy catalog, backend liveness và catalog trả200; catalog có9 sản phẩm. Đây là kiểm tra khả năng truy cập và dữ liệu trả về, không phải luồng đặt hàng E2E. \\
\bottomrule\end{longtable}

Các kết quả trên được tổng hợp từ báo cáo kiểm thử và dữ liệu CI của dự án,
không phải các lượt chạy mới trong quá trình biên tập tài liệu. Phạm vi unit,
integration, HTTP E2E và Sandbox được phân biệt để tránh cộng trùng hoặc suy
rộng kết quả. Các hạn chế còn ảnh hưởng đến kết luận được tổng hợp ở
mục\ref{sec:limitations}.

\section{Triển khai, vận hành và bảo trì}
\label{sec:deployment-release-delivery}
\subsection{Mô hình triển khai và khả năng truy cập}
Frontend được cung cấp tại \url{https://brewlite-a2-three.vercel.app/}; backend
tại \url{https://brewlite-a2-api-hntn.onrender.com/}. Next.js trên Vercel chuyển
yêu cầu qua Route Handler tới NestJS trên Render. PostgreSQL được quản lý
trên Render; mô hình mục tiêu sử dụng kết nối nội bộ giữa API và database.

Ảnh triển khai ngày13/11/2026 ghi Render Auto-Deploy thành công, source
\texttt{02452f228712c3501c8d27f0f1ae737ad206785d}, thời điểm15:39:19 GMT+7.
Dashboard database tại mốc này ghi PostgreSQL18, Singapore, dung lượng1GB,
gói Free và hạn sử dụng10/12/2026. Đây là thông tin tại thời điểm ghi nhận,
không phải cam kết dịch vụ liên tục.

Kiểm tra HTTP ngày14/11/2026 lúc10:31 UTC cho kết quả frontend, proxy
\path{/api/backend/products}, backend \path{/} và \path{/products} đều trả200;
catalog trả JSON gồm9 sản phẩm. Một lượt proxy trước đó bị timeout rồi lần
kiểm tra sau thành công; chưa xác định nguyên nhân timeout. \path{GET /}
trả Hello World! là liveness, không truy vấn database. \path{GET /products}
đọc catalog qua ProductsService và Prisma. Hai phép kiểm tra có mục đích
khác nhau và không thay kiểm thử hành trình trình duyệt.

\subsection{Cấu hình và tích hợp liên tục}
Ứng dụng tổ chức thành \path{apps/web} và \path{apps/api}, mỗi app có lockfile
và scripts riêng. CI dùng Node24; Docker dùng Node26. Pipeline cài dependency
bằng npm ci, lint, unit test và build; API generate Prisma trước khi chạy.
Render có cấu hình Auto-Deploy sau CI theo hồ sơ triển khai. Việc triển khai
cần lưu phiên bản ứng dụng và migration tương ứng để hỗ trợ truy vết sự cố.

Secrets như DATABASE\_URL, JWT signing key và VNPay HashSecret được quản lý
tại backend/platform, không lưu trong repository hoặc đưa vào frontend.
Tên BACKEND\_INTERNAL\_URL chỉ là tên biến cấu hình proxy; không chứng minh
đường truyền giữa Vercel và Render là mạng riêng.

\subsection{Migration và dữ liệu khởi tạo}
Source có10 migrations. Script db:migrate thực thi Prisma migrate deploy,
áp dụng migration có sẵn, có thể thay cấu trúc database nhưng không sinh
file migration mới. Docker target migrate chứa CLI, schema và config;
runtime image chứa dist và dist-seed, không chứa Prisma CLI. Runner migration
phải có kết nối database thích hợp và ghi log status--deploy--status.

Seed local dùng database demo, upsert catalog9 sản phẩm cùng metadata/giá,
giữ stock và không xóa order. Seed Render riêng mặc định dry-run, kiểm tra
đúng target và yêu cầu bước Apply có kiểm soát. Fixture gồm4 categories,
5 toppings,9 products,13 product--topping links và5 coupons. Marker
\texttt{render-demo-v1} trong seed\_executions được ghi cùng transaction;
nếu marker tồn tại mà dữ liệu sai lệch, quy trình dừng để đối chiếu thay vì
xóa marker và chạy lại tùy tiện.

\subsection{Sao lưu, phục hồi và rollback}
Rollback ứng dụng cần một artifact cũ còn khả dụng và tương thích API,
schema, biến môi trường. Vercel Instant Rollback phụ thuộc deployment đủ
điều kiện; Render rollback phụ thuộc artifact được giữ lại và cấu hình
Auto-Deploy \cite{renderdeploy,vercelrollback}. Sau rollback cần kiểm tra
liveness, catalog và các luồng nghiệp vụ bị tác động.

Rollback code không đảo migration hoặc phục hồi dữ liệu. Phục hồi database
cần backup, đích cô lập và kiểm tra dữ liệu trước khi chuyển kết nối.
Runbook mô tả pg\_dump/pg\_restore phù hợp PostgreSQL18, quản lý credential
qua môi trường, tách backup đã tạo với backup đã thử phục hồi. Với database
có thời hạn, dữ liệu cần được bảo toàn hoặc chuyển sang tài nguyên phù hợp
trước ngày hết hạn; hướng dẫn không đồng nghĩa thao tác đó đã được thực hiện.

\subsection{Bảo trì và bàn giao}
Theo dõi liveness, catalog, lỗi API, expiry và reconciliation sau restart
giúp phát hiện lỗi vận hành. Scheduler chạy trong NestJS process nên chịu
ảnh hưởng khi service sleep hoặc ngừng hoạt động. Thay đổi cần kiểm tra
tương thích API/schema, chạy CI và cập nhật runbook. Quyền vận hành được
bàn giao qua membership của nền tảng; secrets dùng cơ chế quản lý riêng.
